% ============================================================
%  example_medium.tex  ——  进阶示例：各种边界写法
%  用法：math_guard.exe example_medium.tex
%
%  这个文件混了三类：
%    【真错误】必须报
%    【正确写法】工具应放过（这些在真实论文里大量出现）
%    【无法判定】工具应承认没判定
% ============================================================
\documentclass{article}
\usepackage{amsmath,amssymb}
\begin{document}

\section{真错误（应报）}

% 符号算错：(-3)×(-5) 应为 +15
\[ (-3)\times(-5)=-15 \]

% 阶乘算错：5! = 120，不是 119
\[ 5!=119 \]

% 根号算错：\sqrt{16} = 4，不是 5
\[ \sqrt{16}=5 \]

% 环境不匹配：\begin{align} 配 \end{gather}
\begin{align}
  x &= 1
\end{gather}

% 未定义引用
See equation \eqref{eq:does-not-exist}.

\section{正确写法（工具应放过，不要报）}

% 勒让德符号（不是除法）
\[ \left(\frac{5}{11}\right)=1 \]

% 有限群里的逆元（不是 1/9）
\[ 9=5^{-1} \]

% 阶乘（!= 不是「不等于」）
\[ 6\cdot 3!=36 \]

% 上标里的装饰记号（⟨1⟩ 是标记，不是指数）
\[ 4^{\langle1\rangle}=10 \]

% 求值竖线（条件等号不是数值断言）
\[ t|_{n=0}=1 \]

% 序理论签名
\[ F=\langle W,\leq\rangle,\qquad L=\{<,c_n:n\geq1\} \]

% 商集
\[ c\in Z/_{\sim}=Q(Z) \]

% TikZ 风格的空占位
\[ x = {} \]

% 中文句中的符号代用
题设常数 $=2\times10^{-4}$，误差 $\le5\%$。

\section{无法判定（工具应承认）}

% 纯符号等式：不代值无法判定
\[ a^{2}+b^{2}=c^{2} \]

% 含变量：需要代入具体数值
\[ \Gamma_{0}/(1-\Gamma_{0})=z(s+z)/(1-z^{2}) \]

% 矩阵
\[ \begin{pmatrix} 1 & 2 \\ 3 & 4 \end{pmatrix} \]

% 多层求和
\[ \sum_{i=1}^{n}\sum_{j=1}^{m} a_{ij} \]

\end{document}
